
	
\chapter{Bilancio di Energia e Sbilancio di Entropia}

Oltre alle leggi di bilancio delle forze, la risposta di un corpo è governata da  
\begin{itemize}
	\item[(i)] la legge di bilancio dell'energia e  
	\item[(ii)] una legge di sbilancio dell'entropia.
\end{itemize}

Queste due leggi sono note come  \textit{prima} e  \textit{seconda legge della termodinamica}. Come già fatto nel caso del bilancio di forze e momenti, formuleremo inizialmente queste due leggi in modo globale per le parti \( \Pi \), e successivamente, utilizzando il requisito che le parti sottostanti siano arbitrarie e il lemma di localizzazione, deriveremo forme locali delle leggi sotto forma di equazioni differenziali alle derivate parziali.


\section{Bilancio di Energia: Prima Legge della Termodinamica}
Definiamo  \textit{moto}  una famiglia mono-parametrica di deformazioni:
	\[
	\mathcal{M} = \left\{ \fb_t \mid t \in T \subset \mathbb{R} \right\}, \qquad 	(\xb, t) \mapsto \yb = \fb(\xb, t) = \fb_t(x);
	\]
il parametro $t$ indica tipicamente il \textit{tempo}. In un moto possiamo definire la \textit{velocità}

	\[
	\mathbf{v}(\xb) = \partial_t \yb = \partial_t \fb(\xb, t)= \partial_t \ub(\xb, t).
	\]

Un coppia costituita da \( \Mc \) e da un campo tensoriale simmetrico \( \Tb \), con \( T(\cdot, t) \) liscio per ogni \( t \in T \), è chiamata un \textit{processo dinamico}.

In aggiunta alla descrizione del moto e all'associato sforzo, saremo interessati a descrivere l'evoluzione di un sistema materiale in termini dei suoi stati termici e meccanici, analizzando dunque un \textit{processo termodinamico}, in cui considereremo non solo sforzo e deformazione, ma anche quantità come temperatura, energia interna, entropia, flusso di calore.

Vediamo adesso quali sono le leggi che governano i processi termodinamici.

La \textit{prima legge della termodinamica} rappresenta un bilancio tra il tasso di variazione dell'energia interna e il tasso al quale l'energia  è trasferita, sotto forma di calore, ad una parte  \( \Pi \) più la potenza meccanica spesa. Nel seguito ci limiteremo a considerare piccole deformazioni.

Introduciamo l'\textit{energia interna} $E(\Pi)$ della parte $\Pi$, espressa in termini di una densità di energia interna per unità di volume $\varepsilon$:
\begin{equation}
	E(\Pi) = \int_{\Pi} \varepsilon.
\end{equation}

La sua derivata temporale, che indichiamo con $\dot{E}(\Pi)$ rappresenta il tasso di aggiunta di energia alla parte $\Pi$.
Consideriamo due tipi di apporto energetico:

\begin{enumerate}
	\item \textit{Sorgenti Meccaniche}. Il sistema di forze  che agisce su $\Pi$, in virtù del fatto che i punti di $\Pi$ assumono una verta velocità durante il moto, spende potenza. In particolare, definiamo la \textit{potenza delle forze} agenti su $\Pi$, o \textit{potenza esterna} la quantità
	\begin{equation}\label{Wext}
		\Wc^{\rm ext}(\Pi)= \int_{\Pi} \bb \cdot \vb + \int_{\partial \Pi} \tb\cdot \vb.
	\end{equation}

Il primo contributo è la potenza spesa dalle forze di volume, mentre il secondo la potenza spesa dalle tensioni che agiscono sulla frontiera di $\Pi$, che adesso abbiamo assunto totalmente interna a $\Omega$.

Usando il teorema di Cauchy, possiamo scrivere che 
	
\begin{equation}
		\Wc^{\rm ext}(\Pi)=	\int_{\Pi} \bb \cdot \vb + \int_{\partial \Pi} \Tb \nb\cdot \vb = \int_{\Pi} \bb \cdot \vb + \int_{\Pi} \Div(\Tb^\top\vb);
	\end{equation}
usando l'identità \eqref{identitadiff}$_5$ e rammentando che $\Tb\in\Sym$, otteniamo che 
\begin{equation}
		\Wc^{\rm ext}(\Pi)	= \int_{\Pi} (\Div\Tb+\bb) \cdot \vb +  \Tb \cdot \dot{\Eb},
\end{equation}
dove abbiamo indicato con
\begin{equation}
\dot\Eb=\sym \nabla \dot \ub=\sym \nabla \vb
\end{equation}
la \textit{velocità di deformazione}. Ricordando che, per l'equazione puntuale di equilibrio $\Div\Tb+\bb=\mathbf{0}$, concludiamo che
\begin{equation}
	\Wc^{\rm ext}(\Pi)= \int_{\Pi}  \Tb \cdot \dot{\Eb}=	\Wc^{\rm int}(\Pi),
\end{equation}
dove $\Wc^{\rm int}(\Pi)$ indica la \textit{potenza interna}, o \textit{potenza dello sforzo}. Abbiamo quindi concluso che la potenza esterna eguaglia la potenza dello sforzo.

\item \textit{Sorgenti termiche}. Indichiamo con $\Qc(\Pi)$ il\textit{ trasferimento di calore} alla parte $\Pi$, che consta di due contributi: uno dato dal \textit{flusso di calore} $\qb$ trasferito a $\Pi$ tramite la sua frontiera $\partial\Pi$ e una \textit{sorgente interna} di calore $r$ (generata, ad esempio, da reazioni chimiche, radiazioni, etc):
\begin{equation}
	\Qc(\Pi) = -\int_{\partial \Pi} \qb\cdot \mathbf{n} + \int_{\Pi} r.
\end{equation}
Il segno negativo nel primo contributo è dovuto al fatto che, essendo $\nb$ la normale uscente a $\partial\Pi$, richiediamo che il trasferimento di calore sia non negativo quando il flusso è entrante in $\partial\Pi$. 

\begin{remark}
Le dimensioni fisiche del flusso di calore e della sorgente interna di calore sono rispettivamente
\begin{equation}
[\qb]=\left[\frac{{\rm energia}}{{\rm tempo}\times{\rm area}} \right], \qquad [r]=\left[\frac{{\rm energia}}{{\rm tempo}\times{\rm volume}} \right]
\end{equation}
\end{remark}


Utilizzando il teorema della divergenza, il trasferimento di calore $\Qc (\Pi)$ può essere riscritto come
\begin{equation}
\Qc(\Pi)= \int_{\Pi}(-\Div \qb +r).
\end{equation}


\end{enumerate}


Possiamo a questo punto stabilire che il bilancio di energia corrisponde alla seguente richiesta:
\begin{equation}
\dot{E}(\Pi)=\Wc^{\rm int}(\Pi)+\Qc (\Pi),
\end{equation}
ovvero
\begin{equation}
\left(  \int_{\Pi}\varepsilon\right)^{\mathop{\vcenter{\hbox{\Large$\cdot$}}}}= \int_{\Pi} \Tb \cdot \dot{\Eb} - \Div \, \mathbf{q} + r
\end{equation}
Poiché $\Pi$ non dipende dal tempo, possiamo quindi stabilire che 
\begin{equation}
\int_{\Pi}\dot{\varepsilon}= \int_{\Pi}  \Tb \cdot \dot{\Eb} - \Div \, \mathbf{q} + r.
\end{equation}
Utilizzando il solito lemma di localizzazione otteniamo allora la seguente

\begin{definition}
\textit{\textbf{Equazione locale di bilancio dell'energia}}:
\begin{equation}\label{bilen}
\dot{\varepsilon} =\Tb \cdot \dot{\Eb} - \Div \, \mathbf{q} + r.
\end{equation}
\end{definition}


\section{Sbilancio di Entropia: Seconda Legge della Termodinamica}
La spesa di potenza rappresenta un trasferimento macroscopico di energia, poiché è calcolata utilizzando la velocità dei punti materiali. Consideriamo il calore come una forma aggiuntiva di trasferimento di energia dovuta alle \textit{fluttuazioni} di atomi e/o molecole, e l'\textit{entropia} come una misura del disordine nel sistema indotto da queste fluttuazioni: maggiore è il grado di disordine, maggiore è l'entropia. Come per l'energia, regioni del corpo  possiedono entropia, e l'entropia può fluire da una regione all'altra e anche da una regione verso il mondo esterno. Ma, diversamente dall'energia, alle parti del corpo è consentito \textit{produrre} entropia.

Introduciamo la \textit{produzione interna di entropia} della parte $\Pi$
\begin{equation}
\Delta (\Pi)=\int_{\Pi}\delta,
\end{equation}
dove $\delta$ ne rappresenta la densità. Ammettiamo che questa produzione  abbia due contributi:
\begin{equation}
\Delta(\Pi)= \dot{\Hc}(\Pi)-\Dc(\Pi);
\end{equation}
$\dot{\Hc}(\Pi)$ rappresenta il tasso a cui l'entropia di $\Pi$ cresce; $\Dc(\Pi)$ è invece il tasso a cui l'entropia è trasferita a $\Pi$.

La premessa di base che i sistemi tendano ad aumentare il loro grado di disordine si manifesta nella richiesta che la produzione interna di entropia sia non negativa in ogni regione $\Pi$:
\begin{equation}
	\Delta(\Pi)\geq 0.
\end{equation}

Come per l'energia, assumiamo che esista un campo scalare $\eta$, la \textit{densità di entropia} tale che 
\begin{equation}
	\Hc(\Pi)=\int_{\Pi}\eta.
\end{equation} 
In analogia con il trasferimento di calore, supponiamo che l'apporto di entropia sia caratterizzato da un flusso $\hb$ e una sorgente $s$:
\begin{equation}
	\Dc(\Pi)=-\int_{\partial\Pi}\hb \cdot \nb +\int_{\Pi} s,
\end{equation}
da cui otteniamo
\begin{equation}\label{sbentr}
\left( \int_{\Pi}\eta \right)^{\mathop{\vcenter{\hbox{\Large$\cdot$}}}}\geq -\int_{\partial\Pi}\hb \cdot \nb +\int_{\Pi} s,
\end{equation}
che rappresenta la forma che diamo alla seconda legge della termodinamica, nel presente contesto.

\begin{remark}
La teoria che stiamo costruendo è coerente con due celebri affermazioni di \textsc{Rudolph Clausius}:
\begin{itemize}
	\item \textit{Die Energie des Weltalls ist konstant }(l'energia dell'universo è costante).
	\item \textit{Die Entropie des Weltalls strebt einem Maximum zu}  (l'entropia dell'universo tende ad un massimo).
\end{itemize}
La parola \textit{Weltall} può essere più prosaicamente sostituita da sistema isolato, ovvero un sistema per cui $\qb\cdot \nb=\hb\cdot \nb=0$ e $r=s=0$; in questo caso, il bilancio di energia e lo sbilancio di entropia si riducono a 
\begin{equation}
\left(\int_{\Pi}\varepsilon\right)^{\mathop{\vcenter{\hbox{\Large$\cdot$}}}}=0, \qquad \left( \int_{\Pi}\eta \right)^{\mathop{\vcenter{\hbox{\Large$\cdot$}}}}\geq 0,
\end{equation}
che corrospondono alle due frasi di Clausius.
\end{remark}

Un'ipotesi fondamentale della teoria relaziona il flusso di entropia e di calore e le due sorgenti; in particolare si assume che esista un campo scalare $\vartheta>0$, la \textit{temperatura assoluta}, tale
 che
 \begin{equation}
 \hb =\frac{\qb}{\vartheta}, \qquad s=\frac{r}{\vartheta}.
 \end{equation}
 
 Flusso di entropia e flusso di calore hanno la stessa direzione e non può annullarsi l'uno senza l'altro. L'ipotesi riflette il fatto che un flusso di calore ha differenti effetti entropici a seconda della temperatura in cui avviene; in particolare, il calore ricevuto a basse temperature è più ``efficace'' di quello ricevuto ad alte temperature.
 
 Possiamo quindi scrivere lo sbilancio di entropia \eqref{sbentr} nella forma della
 
 \begin{definition}
 \noindent \textbf{\textit{Disuguaglianza di Clausius-Duhem}}:
 \begin{equation}
 	\left( \int_{\Pi}\eta \right)^{\mathop{\vcenter{\hbox{\Large$\cdot$}}}}\geq -\int_{\partial\Pi}\frac{\qb}{\vartheta} \cdot \nb +\int_{\Pi} \frac{s}{\vartheta}.
 \end{equation}
 \end{definition}
 
 Utilizzando il solito lemma di localizzazione perveniamo alla 
 \begin{definition}
\textbf{\textit{Disequazione (locale) di sbilancio di entropia}}:
\begin{equation}
	\dot \eta \geq -\Div \left( \frac{\qb}{\vartheta} \right)+\frac{r}{\vartheta}
\end{equation}
 \end{definition}
 e all'espressione della densità di produzione interna di entropia
 \begin{equation}\label{delta}
 \delta=\dot{\eta}+\Div \left( \frac{\qb}{\vartheta} \right)-\frac{r}{\vartheta}\geq 0.
 \end{equation}
 
 Quando all'interno di un sistema avvengono processi irreversibili, come la conduzione di calore, si ha una produzione interna di entropia $\delta$, che misura il tasso di entropia generata per unità di volume a causa di tali processi irreversibili. Abbiamo visto come la temperatura assoluta $\vartheta$ serva come fattore che scala la conversione della produzione di entropia in energia dissipata. La quantità
\begin{equation}
\gamma\coloneqq\delta\vartheta
\end{equation}
rappresenta il \textit{tasso di dissipazione di energia} (o \textit{potenza dissipata}) per unità di volume nel sistema. Infatti, in un sistema termodinamico, il secondo principio della termodinamica stabilisce che la produzione di entropia è sempre associata a una perdita di energia disponibile: questa perdita si manifesta come dissipazione di energia interna.

Utilizzando la \eqref{delta}, possiamo scrivere la potenza dissipata come
\begin{equation}
	\gamma=\vartheta\dot\eta+\vartheta\Div \left( \frac{\qb}{\vartheta} \right)-r\geq 0.
\end{equation}

Utilizzando l'identità \eqref{identitadiff}$_3$, possiamo dedurre che
\begin{equation}
\Div \left( \frac{\qb}{\vartheta} \right)=\nabla\left(\frac{1}{\vartheta}\right)\cdot\qb +\frac{1}{\vartheta}\Div \qb =-\frac{1}{\vartheta^2}\nabla\vartheta\cdot \qb +\frac{1}{\vartheta}\Div \qb 
\end{equation}
 e quindi
 \begin{equation}\label{disga}
	\gamma=\vartheta\dot\eta-\frac{1}{\vartheta}\nabla\vartheta\cdot \qb +\Div \qb -r\geq 0.
 \end{equation}
 
 Utilizzando l'equazione locale di bilancio dell'energia \eqref{bilen} possiamo stabilire allora che
\begin{equation}
	\gamma=\vartheta\dot\eta-\frac{1}{\vartheta}\nabla\vartheta\cdot \qb +\Tb \cdot \dot{\Eb}-\dot{\varepsilon}\geq 0. 
\end{equation}
Se la quantità $\gamma=\delta\vartheta$ rappresenta potenza dissipata per unità di volume, ossia la densità di energia che viene convertita in calore e aumenta l'entropia del sistema a causa di processi irreversibili, possiamo interpretare la quantità $\vartheta\eta$ come la densità di energia  ``intrappolata'' nell'entropia del sistema. 
La differenza tra la densità di energia totale $\varepsilon$ e quella ``intrappolata'' nell'entropia rappresenta l'energia a disposizione del sistema. Introduciamo dunque l'\textit{energia libera di Helmholtz} per unità di volume
\begin{equation}
\psi=\varepsilon-\vartheta\eta.
\end{equation}
 
 La \eqref{disga} può allora essere riscritta come la
 \begin{definition}
\textbf{\textit{Disuguaglianza di dissipazione ridotta}}:
\begin{equation}\label{DDR}
-\gamma=\dot{\psi}+\eta \dot{\vartheta}-\Tb\cdot\dot{\Eb}+\frac{1}{\vartheta}\qb \cdot \nabla \vartheta\leq 0.
\end{equation}
\end{definition} 
I termini della \eqref{DDR} hanno le dimensioni di (energia)/tempo. Integrando sulla parte $\Pi$ possiamo interpretare i vari termini:
\begin{equation}
\underbrace{\int_{\Pi}\gamma}_{{\rm dissipazione}\geq 0}=\underbrace{\int_{\partial \Pi}\Tb \nb \cdot \vb +\int_{\Pi}\bb\cdot \vb}_{{\rm potenza\; esterna}}-
\underbrace{\left(\int_{\Pi}\psi\right)^{\mathop{\vcenter{\hbox{\Large$\cdot$}}}}}_{\scriptsize \shortstack{\text{tasso di variazione} \\ \text{dell'energia libera}}}
-\underbrace{\int_{\Pi}\eta \dot{\vartheta}+\frac{1}{\vartheta}\qb \cdot \nabla\vartheta}_{\text{energia termica}}
\end{equation}

Dettagliamo meglio il significato dell'ultimo contributo, l'energia termica. L'entropia $\eta$ è associata all'energia ``immagazzinata'' termicamente nel sistema. A temperature più elevate, ogni unità di entropia corrisponde a una quantità maggiore di energia termica. Quando la temperatura varia con il tempo $\dot{\vartheta}$, anche l'energia termica associata a questa entropia cambia. In altre parole, il termine $\eta\dot{\vartheta}$ rappresenta l'energia termica per unità di volume che viene aggiunta o sottratta dal sistema per via del cambiamento della temperatura. Se $\dot{\vartheta}$ è positivo, la temperatura aumenta, e quindi il sistema guadagna energia termica; se $\dot{\vartheta}$ è negativo, il sistema perde energia termica.

Il secondo termine  riguarda il flusso di calore $\qb$ e il gradiente di temperatura $\nabla\vartheta$. La quantità $\frac{1}{\vartheta}\qb \cdot \nabla\vartheta$ rappresenta la dissipazione interna associata alla conduzione del calore. In particolare, $\qb \cdot \nabla\vartheta$ descrive quanto calore fluisce attraverso il sistema in presenza di un gradiente di temperatura; dividendo per $\vartheta$, otteniamo una quantità proporzionale alla generazione o assorbimento di entropia associato a questo flusso di calore, che si riflette nel cambiamento dell'energia termica del sistema.


\begin{remark}\label{bb}
Le forze di volume $\bb$ possono essere decomposte in una parte inerziale ed una non inerziale 
\begin{equation}
\bb =\bb^{({\rm in})}+\bb^{({\rm ni})},
\end{equation}
essendo le forze inerziali legate all'accelerazione e alla densità di massa $\rho$:
\begin{equation}
\bb^{({\rm in})}=-\rho \ddot{\ub}=-\rho \dot{\vb}.
\end{equation}
Se si fa questo, l'equazione di bilancio delle forze diventa un'\textit{equazione di evoluzione}, che regola la dinamica dei corpi continui, interpretabile come un bilancio della quantità di moto,
\begin{equation}
\Div \Tb+\bb^{({\rm ni})}=\dot{\pb},
\end{equation}
essendo $\pb\coloneqq\rho \vb$ la \textit{quantità di moto}. La potenza esterna può allora essere valutata come segue:
\begin{equation}
	\Wc^{\rm ext}(\Pi)= \int_{\Pi} \bb^{({\rm ni})} \cdot \vb-\rho \dot{\vb} \cdot \vb + \int_{\partial \Pi} \Tb\nb\cdot \vb=\int_{\Pi} \bb^{({\rm ni})} \cdot \vb-\pb \cdot \dot{\vb} + \int_{\partial \Pi} \Tb\nb\cdot \vb.
\end{equation}
Il termine $-\pb \cdot \dot{\vb}$ ha un'analogia formale e fisica con il termine $-\eta \dot{\vartheta}$ presente nell'energia termica: così come la quantità di moto più essere pensata come una misura della \textit{riluttanza alla quiete} della materia, così l'entropia può essere pensata come una misura della \textit{riluttanza all'ordine}.
\end{remark}

	
	